% ===================================================================
%  TABELO N° 8 — La rekolto.
%  Feuillets 57 a 62 du fac-simile = folios 55 a 60.
%  Correspondance : folio imprime = numero de feuillet - 2.
%
%  Ca tabelo kombinas du « ceno » (scenes) : la Unesma ceno (La
%  Rekolto : la aspekti di la ruro, la pastoro, la rekolto propra, la
%  sturmo) e la Duesma ceno (La Chasado, la Vit-rekolto, Biro e Cidro,
%  la Peskado). La numerado dil alinei e dil renvoyi (1), (2)... esas
%  propra a singla ceno, e riprenas de 1 che la Duesma ceno — d'ube
%  la repeto dil clefi %%K t08-01, t08-02... por la duesma ceno : lo
%  esas l'imprimo mem qua restartigas, ne la transskripto.
% ===================================================================

% --- Feuillet 57 (folio 55) -----------------------------------------
\begin{VUpage}[57]{55}
\VUnotes{143.2mm}{%
\textit{(*) Pro ke ta vorto esas ne preciza : kad esas lin-rekolto,}\nl
\textit{vit-rekolto, pom-rekolto ? Forsan esus bona uzar «} \VUgras{mesono} \textit{»}\nl
\textit{propozita en Mondo XIV, p. 242. Ma til nun ta nova radiko}\nl
\textit{ne esas adoptita dal Akademio e vartas diskuti segun la}\nl
\textit{kustumata regularo Idistala.}%
}
%%K t08-apar-1 apar
\VUsaut{5.0mm}
\VUtitre{15.0pt}{60}{TABELO N\textsuperscript{o} 8}
\VUsaut{3.0mm}
\VUcentre{13.2pt}{90}{\textit{Unesma ceno.}}
\VUsaut{3.0mm}
\VUpk{10.2pt}{40}{La Rekolto (*).}
\VUsaut{4.0mm}

%%K t08-tit-1 sub
\VUpk{10.2pt}{40}{La Aspekti di la Ruro.}
\VUsaut{3.0mm}

%%K t08-c1-01-1 p
\VUblancAlinea
1. --- Kun la \VUgras{somero}, unplusafoye chanjis\nl
l'aspekto di la ruro. La semino konfidita al\nl
sulko jermifis; verda planteto ekiris de la sulo\nl
e spikifis. La grano esis formacata, pose lu\nl
matureskis e mem de nun on bezonas nur re\cc
koltar.

%%K t08-c1-02-1 p
\VUblancAlinea
2. --- La \VUgras{rural flori} esas desklozita. \VUgras{Aciani} \textsuperscript{(1)},\nl
\VUgras{papavereti} \textsuperscript{(2)} e \VUgras{digitali} \textsuperscript{(3)} mixas al uniforma\nl
nuanco dil frument-agri la gaya kontrasto di\nl
sua fresha \VUgras{koroli.}

%%K t08-c1-03-1 p
\VUblancAlinea
3. --- La frukt-arbori su garnisis per folii;\nl
la frukto matureskis. La mikra \VUgras{ceriziero} \textsuperscript{(4)}\nl
abundas de bela \VUgras{cerizi} \textsuperscript{(5)} quin la pueri ple\cc
zurege kolias e manjas.

%%K t08-c1-04-1 p
\VUblancAlinea
4. --- De nun la brutaro povas pasar la tota\nl
jorno en la liber aero.

%%K t08-c1-04-2 p
\VUblancAlinea
La \VUgras{mutonyuni} \textsuperscript{(6)} grandeskis e divenis bel\nl
\VUgras{mutonini} \textsuperscript{(7)} e bel \VUgras{mutonuli} \textsuperscript{(8)} kun densa lan\cc
felo. Li tranquile pasturas la \VUgras{herbo} dil \VUgras{pastu}\cc
\VUgras{reyo} \textsuperscript{(9)} buntizita per \VUgras{margriteti.}
\end{VUpage}

% --- Feuillet 58 (folio 56) -----------------------------------------
\begin{VUpage}[58]{56}
%%K t08-c1-04-3 p
\VUblancAlinea
Balde on tondeskos \VUgras{ta} bestii por havar lia\nl
\VUgras{lano} ek qua on fabrikas la drapo di nia vesti.

%%K t08-tit-2 sub
\VUsaut{5.0mm}
\VUpk{10.2pt}{40}{La Pastoro.}
\VUsaut{3.4mm}

%%K t08-c1-05-1 p
\VUblancAlinea
5. --- La \VUgras{pastoro} \textsuperscript{(10)} semblas ne tre atencar\nl
li. Il sucias nur la sturmo qua pasas ye la ho\cc
rizonto, e la \VUgras{muton-gardisto} \textsuperscript{(13)}, qua proximi\cc
gas \VUgras{gurdo} \textsuperscript{(14)} a sua labii semblas mem plu sen\cc
sucia.

%%K t08-c1-06-1 p
\VUblancAlinea
6. --- La varmeso esas opresanta; nam la\nl
\VUgras{hundo} \textsuperscript{(15)} venas, lu anke, sendurstigar su; lu\nl
lapas la fresh aquo di la \VUgras{rivereto} \textsuperscript{(16)} e pavo\cc
rigas kompatinda \VUgras{raneto} \textsuperscript{(17)} qua blotisabis en\nl
la rivo-herbi. Ica saltas aden la klar aquo ube\nl
florifas la \VUgras{nenufari} \textsuperscript{(18)}.

%%K t08-c1-07-1 p
\VUblancAlinea
7. --- \VUgras{Motacilo} \textsuperscript{(19)} su balnas apud la \VUgras{fon}\cc
\VUgras{teto} \textsuperscript{(20)} qua spricas inter du rokaji e su celas\nl
sub \VUgras{l'arbusti dornoza} \textsuperscript{(21)} e la \VUgras{krateg-bushi} \textsuperscript{(22)}.

%%K t08-tit-3 sub
\VUsaut{5.0mm}
\VUpk{10.2pt}{40}{La Rekolto.}
\VUsaut{3.4mm}

%%K t08-c1-08-1 p
\VUblancAlinea
8. --- Por la rurana pueri, la frumento-re\cc
kolto esas tre amuzant epoko. Li joyoze kul\cc
\VUgras{butas}\textsuperscript{(24)} sur la \VUgras{palio-amasi} \textsuperscript{(23)}, li helpas la\nl
\VUgras{rekoltisti} \textsuperscript{(26)} e dum ke ici falchas la \VUgras{frument}\cc
\VUgras{agro} \textsuperscript{(27)} per sua \VUgras{falchili} \textsuperscript{(25)} bone akutigita per\nl
la \VUgras{grindo-petro} \textsuperscript{(30)}, li kolektas la frumento e\nl
ligas olu kom \VUgras{garbi} \textsuperscript{(31)} o kom \VUgras{fasketi} \textsuperscript{(32)}.

%%K t08-c1-09-1 p
\VUblancAlinea
9. --- Kelkafoye on permisas a le maxim\nl
granda inter li sekar per \VUgras{siklo} \textsuperscript{(12)} la \VUgras{orea}\nl
\VUgras{stipi} \textsuperscript{(28)} qui portas la grava \VUgras{spiki} \textsuperscript{(29)} grano\cc
plena; e la mikri, surveyante la \VUgras{brutaro} \textsuperscript{(33)}
\parplein
\end{VUpage}

% --- Feuillet 59 (folio 57) -----------------------------------------
\begin{VUpage}[59]{57}
%%K t08-c1-09-1 p suite
\VUcontinue
qua pasturas en la \VUgras{pastureyo} \textsuperscript{(34)} sub la ombro\nl
dil granda \VUgras{tilio} \textsuperscript{(35)} kaptas per sua \VUgras{reto} \textsuperscript{(36)} la\nl
bela \VUgras{papilioni.}

%%K t08-c1-09-2 p
\VUblancAlinea
Mem kelki klimas sur la \VUgras{chareto} \textsuperscript{(37)} por\nl
ordinar la garbi quin on pasigas a li per\nl
\VUgras{forko} \textsuperscript{(38)}.

%%K t08-c1-10-1 p
\VUblancAlinea
10. --- Sur la tota \VUgras{planajo} \textsuperscript{(39)} ube flugeskas\nl
la \VUgras{alaudi} \textsuperscript{(41)} regnas l'animeso. Omni esas oku\cc
pata pri la rekolto. Ma, dum ke la povri duras\nl
uzar la olda instrumenti, \VUgras{falchili, sikli} e \VUgras{dra}\cc
\VUgras{shili} \textsuperscript{(40)}, la richa proprieteri falchigas sua fru\cc
mento o sua \VUgras{sekalo} \textsuperscript{(43)} per \VUgras{falcho-mashino} \textsuperscript{(42)}.\nl
Pose, ne bezonante portar la garbi en la gar\cc
beyo, on facas granda \VUgras{garb-amasi} \textsuperscript{(44)}.

%%K t08-c1-11-1 p
\VUblancAlinea
11. --- Lore on adduktas \VUgras{drashmashino} \textsuperscript{(47)}\nl
quan \VUgras{lokomobilo} \textsuperscript{(45)} movas per \VUgras{transmis-ri}\cc
\VUgras{meno} \textsuperscript{(46)} e tale la grano separesas de la \VUgras{palio,}\nl
sen livar l'agro ube ol esis semata.

%%K t08-tit-4 sub
\VUsaut{5.0mm}
\VUpk{10.2pt}{40}{La Sturmo.}
\VUsaut{3.4mm}

%%K t08-c1-12-1 p
\VUblancAlinea
12. --- Ofte \VUgras{sturmegi} \textsuperscript{(53)} eventas, lor la tempo\nl
dil rekolto. Unesme la grondi dil \VUgras{tondro} esas\nl
audata de fore; \VUgras{ventego de westo} \textsuperscript{(54)} sufles\cc
kas ed anhelante kurvigas la maxim gros ar\cc
bori dil \VUgras{bosko} \textsuperscript{(49)}. Nigra \VUgras{nubegi} \textsuperscript{(56)} asaltas la\nl
cielo; li esas trairata da zigzagi dazlanta dil\nl
\VUgras{fulmino} \textsuperscript{(55)}. La \VUgras{pluvo} e kelkafoye la \VUgras{grelo} fa\cc
leskas, aplastante la rekoltaji pronta por fal\cc
chesor.

%%K t08-c1-13-1 p
\VUblancAlinea
13. --- Ofte la fulmino incendias konstruk\cc
turo. Tale lastayare la tekto dil \VUgras{vento-mue}\cc
\VUgras{lilo} \textsuperscript{(50)} esis cindrigata e du de lua \VUgras{alegi} \textsuperscript{(51)} fra\cc
kasata.
\end{VUpage}

% --- Feuillet 60 (folio 58) -----------------------------------------
\begin{VUpage}[60]{58}
%%K t08-c1-14-1 p
\VUblancAlinea
14. --- Pos kelka hori, la kalmeso rinaskas.\nl
Briloza \VUgras{ciel-arko} \textsuperscript{(52)} aparas en la horizonto e\nl
la naturo riprenas sua vivo interruptita dum\nl
kelka instanti. La rurani, qui refujis en la\nl
\VUgras{kabani} \textsuperscript{(48)} rilaboreskas e kurajoze esforcas re\cc
parar la domaji quin li jus subisis.

%%K t08-tit-5 sub
\VUsaut{6.0mm}
\VUcentre{13.2pt}{90}{\textit{Duesma ceno.}}
\VUsaut{3.6mm}

%%K t08-tit-6 sub
\VUpk{10.2pt}{40}{La Chasado. -- La Vit-rekolto.}
\VUsaut{1.4mm}

%%K t08-tit-7 sub
\VUpk{10.2pt}{40}{La Peskado.}
\VUsaut{4.0mm}

%%K t08-c2-01-1 p
\VUblancAlinea
1. --- Cirkum la fino di \VUgras{agosto} o la mezo di\nl
\VUgras{septembro}, segun la regioni, eventas l'aperto\nl
dil chasado. Apene l'unesma sun-radii igis la\nl
\VUgras{lacerti} \textsuperscript{(2)} ekirar lia truo, kande la \VUgras{chasero} \textsuperscript{(5)}\nl
voyireskas kun sua \VUgras{hundi} \textsuperscript{(4)}.

%%K t08-c2-02-1 p
\VUblancAlinea
2. --- Il metis sua solida \VUgras{chas-kostumo} \textsuperscript{(9)} ek\nl
dika drapo, pedvestizis su per ledra \VUgras{getri,} per\nl
qui il sucesos marchar en la humida herbo ube\nl
celas su la \VUgras{rospi} \textsuperscript{(1)}; il prenis sua \VUgras{fusilo} \textsuperscript{(6)} e\nl
sua \VUgras{vildo-sako} \textsuperscript{(10)}, e yen ke il departas.

%%K t08-c2-03-1 p
\VUblancAlinea
3. --- La fervoro di la persequo adduktas lu\nl
meze di vit-agro ube la vit-rekolto esas tre\nl
vivaca. La pueri dil vit-kultivisto, qui ordinare\nl
surveyas la kaprini sur la voyo, cadie lasas eli\nl
pasturar hazarde, e pos ligir a paliso olda \VUgras{ka}\cc
\VUgras{prulo} \textsuperscript{(3)} qua ne falius profitar sua libereso por\nl
eskapar, li, suafoye, atribuas a su partopreno\nl
en la plezuro. Unu sizas vit-ber-grapo en \VUgras{re}\cc
\VUgras{kolto-korbo} \textsuperscript{(12)} ed amuzas su per fluigar la\nl
\VUgras{suko} \textsuperscript{(14)} en gobleto \textsuperscript{13)} por fabrikar mosto\nl
(vino ne fermentacinta). Altru kuafas su per\nl
\VUgras{folio-krono} \textsuperscript{(15)} quan il jus plektis.
\parplein
\end{VUpage}

% --- Feuillet 61 (folio 59) -----------------------------------------
\begin{VUpage}[61]{59}
%%K t08-c2-03-1 p suite
\VUcontinue
Proxim li, senshama \VUgras{paseri} \textsuperscript{(16)} venas mem\nl
avan la presilo por marodar la vit-beri.

%%K t08-c2-04-1 p
\VUblancAlinea
4. --- Dum ke la pueri su amuzas, la viri\nl
laboras. La \VUgras{vit-rekoltisti} \textsuperscript{(18)} adportas la vit\cc
beri aden la kelero per \VUgras{dorso-korbi} \textsuperscript{(17)}; \VUgras{bare}\cc
\VUgras{lifisto} \textsuperscript{(19)} reparas la \VUgras{bareli} \textsuperscript{(20)}, verifikas ka la\nl
\VUgras{duvi} \textsuperscript{(22)} e la \VUgras{fundi} \textsuperscript{(21)} esas bonstanda, e rem\cc
plasigas la \VUgras{ringegi} \textsuperscript{(23)} ruptita. Lu anke laboris\nl
la \VUgras{kuvegi} \textsuperscript{(25)} en qui on varsas la \VUgras{vitberi} \textsuperscript{(26)} por\nl
igar li fermentacar.

%%K t08-c2-05-1 p
\VUblancAlinea
5. --- Kande la fermentaco parfinis, on rakas\nl
la vino ed on pozas la reziduo sur la \VUgras{presilo} \textsuperscript{(28)}\nl
e progresive klemante per la \VUgras{skrubego} \textsuperscript{(29)} on\nl
ekpresas la suko qua fluas en \VUgras{kuveto} \textsuperscript{(32)} ed\nl
ankore on obtenas \VUgras{vino} \textsuperscript{(31)}.

%%K t08-tit-8 sub
\VUsaut{6.0mm}
\VUpk{10.2pt}{40}{Biro e Cidro.}
\VUsaut{3.4mm}

%%K t08-c2-06-1 p
\VUblancAlinea
6. --- An la muro dil \VUgras{kelero} \textsuperscript{(27)} klimas bran\cc
cho di \VUgras{lupulo} \textsuperscript{(30)}. Ibe ol esas nur orniva planto,\nl
ma en kelka landi, ube la vito esas tre rara,\nl
la lupulo esas kultivata por fabrikar la \VUgras{biro.}

%%K t08-c2-07-1 p
\VUblancAlinea
7. --- La mikra \VUgras{pomiero} \textsuperscript{(33)} esas charjita per\nl
pomi qui produktos kande li esos matura, bo\cc
nega \VUgras{cidro.} En sua \VUgras{kajo} \textsuperscript{(34)} la \VUgras{serino} \textsuperscript{(35)} e la\nl
\VUgras{karduelo} \textsuperscript{(36)} regardas per envidioz okuli la\nl
paseri qui libere petuladas.

%%K t08-c2-08-1 p
\VUblancAlinea
8. --- Til fino di vido, sur la pento dil kolini,\nl
esas aliniita \VUgras{vito-stangi} \textsuperscript{(40)} qui suportas la\nl
belega \VUgras{(vito)-trunki} \textsuperscript{(39)}, garnisita da \VUgras{grapi}\nl
pezoza.

%%K t08-c2-08-2 p
\VUblancAlinea
Dum la tota yaro, la \VUgras{vit-kultivisto} \textsuperscript{(42)} labo\cc
ris por sorgar sua \VUgras{vit-agro} \textsuperscript{(38)} e nun il esas
\parplein
\end{VUpage}

% --- Feuillet 62 (folio 60) -----------------------------------------
\begin{VUpage}[62]{60}
%%K t08-c2-08-2 p suite
\VUcontinue
rekompensata po sua peno per bela rekolto.\nl
La \VUgras{vit-rekoltistini} \textsuperscript{(41)} plenigas la kuvi pozita\nl
sur la chareto tirata da \VUgras{mulini} \textsuperscript{(43)}, ed omni\nl
esas joyoza.

%%K t08-tit-9 sub
\VUsaut{5.0mm}
\VUpk{10.2pt}{40}{La Peskado.}
\VUsaut{3.4mm}

%%K t08-c2-09-1 p
\VUblancAlinea
9. --- Same esas pri la \VUgras{fil-peskeri} \textsuperscript{(44)} qui de\nl
la matino, venis instalar su an la rivo dil aquo\nl
kun sua \VUgras{pesko-bastoni} \textsuperscript{(45)}.

%%K t08-c2-09-2 p
\VUblancAlinea
La \VUgras{fisho} \textsuperscript{(47)} snapas l'\VUgras{angelo} quik kande ili\nl
plunjas sua \VUgras{filo} \textsuperscript{(46)} en l'aquo, ed apene li havis\nl
tempo por rinovigar la \VUgras{lurilo,} kande nova\nl
viktimo friskas ye l'extremajo dil filo.

%%K t08-c2-09-3 p
\VUblancAlinea
Tamen la \VUgras{reto-peskero} \textsuperscript{(48)} qua, de sua \VUgras{bar}\cc
\VUgras{ko} \textsuperscript{(49)}, lansas la \VUgras{(lanso)-reto} aden la mezo dil\nl
\VUgras{rivero} \textsuperscript{(50)} kaptas plu multa fishi.

%%K t08-c2-10-1 p
\VUblancAlinea
10. --- L'autuno esas la sezono dil bona\nl
promenadi : pedire, \VUgras{kalvalke} \textsuperscript{(52)}, \VUgras{bicikle} \textsuperscript{(53)},\nl
\VUgras{automobile} \textsuperscript{(54)}. La \VUgras{voyi} \textsuperscript{(51)} esas neta ed on\nl
profitas la lasta bela dii, nam yen ke ja la\nl
\VUgras{hirundi} \textsuperscript{(56)} asemblas su sur la tekti por fluges\cc
kar per tota ali vers plu varma landi. La fo\cc
liaro dil olda \VUgras{nuciero} \textsuperscript{(57)} flaveskas, la \VUgras{nuci} fa\cc
las e la alt \VUgras{arbori} grupe dispozita quale bu\cc
\VUgras{keto} \textsuperscript{(58)} sur la \VUgras{altajo} \textsuperscript{(59)} balde senfolieskos.

%%K t08-c2-11-1 p
\VUblancAlinea
11. --- La \VUgras{suno} \textsuperscript{(60)} su levas plu tarde, la\nl
jorni deskreskas, koldeteso pasable pikanta\nl
komencas esar sentata, ye la \VUgras{matino,} e yen\nl
ke ja flugi de \VUgras{bekasi} \textsuperscript{(61)} livas nia klimati ne\cc
gastigema.

\VUsaut{6.0mm}
\VUfilet{22mm}
\end{VUpage}
